% \subsection{Hyperparameter Tuning}
%
% AUDA provides two methods for hyperparameter tuning: grid search \cite{hastie2009elements} and random search \cite{bergstra2012random}. Let there be \(d\) hyperparameters with names \(\{ \lambda_1, \lambda_2, \cdots, \lambda_d \}\). For each hyperparameter \(\lambda_j\), denote its domain as \(\mathcal{D}_j = (a_j, b_j)\). The full search space is the Cartesian product \(\mathcal{D} = \mathcal{D}_1 \times \mathcal{D}_2 \times \cdots \times \mathcal{D}_d \).
%
% In random search, for each iteration \(t = 1, 2, \dots, T\), sample a hyperparameter vector \( \bm{\lambda}^{(t)} = (\lambda_1^{(t)}, \dots, \lambda_d^{(t)}) \in \mathcal{D} \). Each hyperparameter can be sampled either uniformaly or log-uniformly. For each dimension \(j = 1, 2, \dots, d\),
%
% \begin{itemize}
% 	\item If \(\lambda_j\) is sampled uniformly, then \(\lambda_j^{(t)} \sim \text{Uniform}(a_j, b_j) \).
% 	\item If \(\lambda_j\) is sampled log-uniformly, then \(\log(\lambda_j^{(t)}) \sim \text{Uniform}(\log a_j, \log b_j)\).
% \end{itemize}
%
% The configuration set \( \mathcal{G} = \{ \bm{\lambda}^{(t)} \mid t = 1, 2, \dots, T \} \).
%
% In grid search, let \( M \in \mathcal{N} \) be the number of grid points for each hyperparameter. For each dimension \(j = 1, 2, \dots, d\), construct a set of grid points \(\mathcal{G}_j\):
%
% \begin{itemize}
% 	\item { For uniform grid:
% 	      \[
% 		      G_j = \left\{ a_j + \frac{m}{M - 1}(b_j - a_j) \middle| m = 0, 1, \dots, M - 1 \right\}
% 	      \]
% 	      }
% 	\item { For log-uniform grid:
% 	      \[
% 		      G_j = \left\{ \exp \left( \log a_j + \frac{m}{M - 1}(\log b_j - \log a_j) \right) \middle| m = 0, 1, \dots, M - 1 \right\}
% 	      \]
% 	      }
% \end{itemize}
%
% Then, configuration set is \(\mathcal{G} = G_1 \times G_2 \times \cdots \times G_d \); each configuration \( \bm{\lambda}^{(t)} \in \mathcal{G} \).
%
% In both cases, for each configuration in \(\mathcal{G}\), the model is trained and evaluated on the validation set using time-series cross-validation \cite{hyndman2018forecasting} for forecasting tasks and repeated random holdout validation \cite{kohavi1995study}. The performance metrics of folds are combined with the configuration to form pairs \( (\bm{\lambda}, (m_1, m_2, \dots, m^K) \), where \(K\) is the number of folds and \(m_k\) is the performance metric on the \(k\)-th fold.
%
% Then, we apply a variant 1-SE rule \cite{little2019statistical} to select the simplest configuration from the candidate configurations. We first compute the mean performance metric across folds for each configuration:
%
% \[
% 	\mu^{(t)} = \frac{1}{K} \sum_{k=1}^K m_k^{(t)}.
% \]
%
% Next, we compute the standard error of the performance metrics for the configuration having the best mean performance \(E = \frac{\sigma}{\sqrt{n}}\). Then, we pick all the configurations whose mean performance is within \(\mu + \alpha E\), where \(\alpha\) is an adjustment coefficient. Finally, we select the simplest configuration among these candidates.
%
% The purpose of the 1-SE rule is to avoid accidental overfitting due to randomness in hyperparameter search. The constant \(\alpha\) controls the trade-off between performance and simplicity. In this study, the iteration of random search is set to \(T = 128\) and the number of grid points for each hyperparameter is set to \(M = 8\). Under these settings, we find that it is proper to set \(\alpha\) to \(0.1\).

\subsection{Hyperparameter Tuning}

AUDA provides two methods for hyperparameter tuning: grid search \cite{hastie2009elements} and random search \cite{bergstra2012random}. Suppose there are \(d\) hyperparameters with names \(\{ \lambda_1, \lambda_2, \dots, \lambda_d \}\). For each hyperparameter \(\lambda_j\), let its domain be \(\mathcal{D}_j = (a_j, b_j)\). The full search space is the Cartesian product
\[
	\mathcal{D} = \mathcal{D}_1 \times \mathcal{D}_2 \times \cdots \times \mathcal{D}_d .
\]

In random search, for each iteration \(t = 1, 2, \dots, T\), a hyperparameter vector
\[
	\bm{\lambda}^{(t)} = \left(\lambda_1^{(t)}, \dots, \lambda_d^{(t)}\right) \in \mathcal{D}
\]
is sampled. Each hyperparameter is sampled either uniformly or log-uniformly. For each dimension \(j = 1, 2, \dots, d\),
\begin{itemize}
	\item if \(\lambda_j\) is sampled uniformly, then
	      \[
		      \lambda_j^{(t)} \sim \mathrm{Uniform}(a_j, b_j);
	      \]
	\item if \(\lambda_j\) is sampled log-uniformly, then
	      \[
		      \log\left(\lambda_j^{(t)}\right) \sim \mathrm{Uniform}(\log a_j, \log b_j).
	      \]
\end{itemize}
The resulting configuration set is
\[
	\mathcal{G} = \left\{ \bm{\lambda}^{(t)} \mid t = 1, 2, \dots, T \right\}.
\]

In grid search, let \(M \in \mathbb{N}\) be the number of grid points for each hyperparameter. For each dimension \(j = 1, 2, \dots, d\), a set of grid points \(\mathcal{G}_j\) is constructed as follows:
\begin{itemize}
	\item for a uniform grid,
	      \[
		      \mathcal{G}_j = \left\{ a_j + \frac{m}{M - 1}(b_j - a_j) \middle| m = 0, 1, \dots, M - 1 \right\};
	      \]
	\item for a log-uniform grid,
	      \[
		      \mathcal{G}_j = \left\{ \exp \left( \log a_j + \frac{m}{M - 1}(\log b_j - \log a_j) \right) \middle| m = 0, 1, \dots, M - 1 \right\}.
	      \]
\end{itemize}
The resulting configuration set is \(\mathcal{G} = \mathcal{G}_1 \times \mathcal{G}_2 \times \cdots \times \mathcal{G}_d\), where each configuration \(\bm{\lambda}^{(t)} \in \mathcal{G}\).

For both tuning methods, each configuration in \(\mathcal{G}\) is used to train the model and evaluate its performance on a validation set. Forecasting tasks use time-series cross-validation \cite{hyndman2018forecasting}, whereas other tasks use repeated random holdout validation \cite{kohavi1995study}. The fold-wise performance metrics are combined with each configuration to form pairs
\[
	\left(\bm{\lambda}^{(t)}, \left(m_1^{(t)}, m_2^{(t)}, \dots, m_K^{(t)}\right)\right),
\]
where \(K\) is the number of folds and \(m_k^{(t)}\) is the performance metric of the \(t\)-th configuration on the \(k\)-th fold.

We then apply a variant of the 1-SE rule \cite{little2019statistical} to select the simplest configuration from the candidate configurations. First, we compute the mean performance metric across folds for each configuration:
\[
	\mu^{(t)} = \frac{1}{K} \sum_{k=1}^K m_k^{(t)}.
\]

Next, we compute the standard error of the performance metrics for the configuration with the best mean performance:
\[
	E = \frac{\sigma}{\sqrt{K}},
\]
where \(\sigma\) is the standard deviation of the fold-wise performance metrics for that configuration. We then identify all configurations whose mean performance is within \(\mu^\star + \alpha E\), where \(\mu^\star\) is the best mean performance and \(\alpha\) is an adjustment coefficient. Finally, we select the simplest configuration among these candidates.

The purpose of the 1-SE rule is to reduce accidental overfitting caused by randomness in hyperparameter search. The constant \(\alpha\) controls the trade-off between performance and simplicity. In this study, the number of random-search iterations is set to \(T = 128\), and the number of grid points for each hyperparameter is set to \(M = 8\). Under these settings, we find that setting \(\alpha = 0.1\) is appropriate.

Table~\ref{tab:complexity-key-hyperparameters} shows the complexity ordering for the hyperparameters of all tunable models. Lower complexity-key values are treated as simpler by the 1-SE rule. For example, in Ridge Regression, a lower degree and a higher alpha contribute to a simpler model.

\begin{table}[t]
	\centering
	\caption{Complexity ordering for hyperparameters of all tunable models.}
	\label{tab:complexity-key-hyperparameters}
	\begin{tabular}{lll}
		\toprule
		Model & Hyperparameter   & Complexity-key contribution \\
		\midrule
		\multirow{2}{*}{Ridge Regression}
		      & degree           & $+\text{degree}$            \\
		      & alpha            & $-\alpha$                   \\
		\midrule
		\multirow{2}{*}{Gaussian Process Regression}
		      & length\_scale    & $-\text{length\_scale}$     \\
		      & noise\_level     & $-\text{noise\_level}$      \\
		\midrule
		\multirow{3}{*}{Support Vector Regression}
		      & C                & $+C$                        \\
		      & epsilon          & $-\epsilon$                 \\
		      & gamma (if tuned) & $+\gamma$                   \\
		\bottomrule
	\end{tabular}
\end{table}
